\section{Construction}\label{sec:agg creds construction}

The protocol involves a registration authority and a collection of issuers.
Each user first establishes a pseudonymous identity with the registration
authority. Issuers then issue credentials to that identity for the attributes
they certify. Before presentation, the user aggregates the credentials
corresponding to the attributes they wish to disclose into a single aggregate
signature. Finally, the user proves knowledge of the aggregate credential and
their registered identity using a signature of knowledge.

Throughout this construction, we refer simply to proofs. These are instantiated
using signatures of knowledge, allowing the resulting proofs to be bound to
protocol messages where required.

\paragraph{Setup.} During setup, the registration authority generates a key
pair for the digital signature scheme and public parameters for the proofs used
throughout the protocol. Each issuer independently generates a key pair for the
aggregate tag based signature scheme. The registration authority's public key,
each issuer's public key, and the public parameters are published.

The registration authority (RA) samples a key pair $(\sk, \pk) \leftarrow
\DSgen(\secparam)$ and publishes~$\pk$. They generates system wide public
parameters $\crs$ for the necessary signatures of knowledge, needed throughout.
Each issuer then generates a key pair $(\isk{i}, \ipk{i}) \leftarrow
\ATSgen(\secparam)$ for the aggregate tag based signature scheme~$\ATS$ and
publishes~$\ipk{i}$.

\paragraph{Registration.} During registration, the user establishes a long term
pseudonymous identity with the registration authority. The user proves
knowledge of the secret key corresponding to their public key, after which the
registration authority binds that public key to a unique user tag.

The user samples $(\sk, \pk) \leftarrow \DSgen(\secparam)$, and sends $\pk$ to
the RA together with a proof on a fresh $\msg$ from the RA that proves the user
knows the corresponding secret key $\sk$.  If the proof verifies, then the RA
assigns the user a unique tag~$\tau$ and issues $\psi \leftarrow \DSsign(\sk,
\pk \| \tau)$.

The user initialises their credential state as $(\tau, \sk, \pk, \psi, \cred_1,
\ldots, \cred_N)$, with $\cred_i = \varnothing$ for all $i \in [N]$.

\paragraph{Issuance.} To obtain an attribute credential, the user authenticates
their registered identity to an issuer. If the issuer is satisfied that the
user possesses the requested attribute value, it issues an aggregate tag-based
signature on that attribute.

On input $(\tau, \pk, \psi, \atv{i})$ from the user, issuer first sets $m =
(\pk, \tau)$, and checks that $\DSverify(\pk, m, \psi) = 1$.  If not, it
rejects.  Otherwise, it verifies that the user knows $\sk$ via a proof.  If
this check succeeds and the user holds attribute value $\atv{i}$, the issuers
sends back $\sigma_i \leftarrow \ATSsign(\isk{i}, \tau, \atv{i})$, and the
user, subsequently, sets $\cred_i = (\atv{i}, \sigma_i)$.

\paragraph{Aggregation.} Before presenting credentials, the user aggregates the
signatures corresponding to the attributes they wish to disclose into a single
aggregate signature.

\paragraph{Presentation.} To prove possession of a collection of attributes,
the user aggregates the corresponding credentials and produces a proof.  This
simultaneously demonstrates possession of a valid registration credential and
that the aggregate signature verifies under the disclosed issuer public keys
and attribute values.

The user has $(\tau, \sk, \pk, \psi, \cred_1, \ldots, \cred_N)$, where $\cred_i
= (\atv{i}, \sigma_i)$.  To prove possession of attribute values $(\atv{i})_{i
\in \ISSUERS}$ to a verifier providing a fresh message $\msg$, the user
computes $\sigma \leftarrow \ATSaggregate(\{\sigma_i\}_{i \in \ISSUERS})$, and
proof $ \proof \leftarrow \soksign(\crs, \inst, \witness, \msg)$, where $\inst
= (\pk, \ISSUERS, \{\ipk{i}, \atv{i}\}_{i \in \ISSUERS})$ and $\witness =
(\tau, \sk, \pk, \psi, \sigma)$, and which proves the following: \begin{align*}
(\sk, \pk) \leftarrow &\DSgen(\secparam) \quad \wedge \quad  \DSverify(\pk, \pk
\|  \tau, \psi) = 1 \\ \wedge &\quad \ATSverifyagg(\{\ipk{i}\}_{i \in
\ISSUERS}, \tau, \{\atv{i}\}_{i \in \ISSUERS}, \sigma) = 1.  \end{align*}

\paragraph{Verification.} The verifier checks the proof with respect to the
disclosed attributes and issuer public keys. The presentation is accepted if
and only if the verification succeeds.

On input $(\{\atv{i}\}_{i \in \ISSUERS}, \msg, \proof)$, the verifier returns
the output of $\sokverify(\crs,\inst,\msg,\proof)$, with $\inst = (\pk,
\ISSUERS, \{\ipk{i}, \atv{i}\}_{i \in \ISSUERS})$.

\begin{theorem}\label{thm:bbsecurity} If $\DS$ is an existentially unforgeable
  digital signature scheme, $\ATS$ an existentially unforgeable aggregate tag
  based signature scheme under tag based aggregation, and $\SOK$ a secure
  signature of knowledge, then the above protocol securely realises ideal
functionality~$\F{\ATS}$ in the standalone model.  \end{theorem}
